\section{问题重述}

体重水平与人体健康状况密切相关，体重异常特别是超重和肥胖是导致心脑血管疾病、糖尿病和部分癌症等慢性病的重要危险因素。调查显示，我国居民超重肥胖形势不容乐观，亟需加强干预，予以改善。2024年6月，国家卫生健康委等16个部门联合制定了《"体重管理年"活动实施方案》，倡导和推进文明健康生活方式，提升全民体重管理意识和技能，预防和控制超重肥胖，切实推动慢性病防治关口前移。

根据中国BMI指数标准，大于或等于24属于超重，大于或等于28属于肥胖。BMI是国际上常用来衡量人体胖瘦程度与健康状态的主要依据，计算方式是体重（公斤数）除以身高（米）的平方。中国目前的BMI标准比国际的更严苛，按照国际标准，介于25至30之间属于超重，大于或等于30属于肥胖。

中国近四年来人口持续下降，超重肥胖率却加速上升，而且年轻化趋势明显，这意味着未来健康劳动力可能减少，加上老龄化加剧，经济增速持续下行，使得应对这一挑战变得更加困难。

经公开来源核验，本文采用5个全国成人超重肥胖率数据点：1992年20.0\%\textsuperscript{[1]}、2002年29.9\%\textsuperscript{[2]}、2012年42.0\%\textsuperscript{[3]}、2020年50.7\%\textsuperscript{[4]}和2023年57.0\%；其中2023年数据以Hu等报告的56.9\%为依据并保留一位小数\textsuperscript{[25]}。

本文需要通过数学建模，求解下列问题：
\begin{enumerate}[label=(\arabic*),leftmargin=2em]
  \item 通过数学建模对我国超重肥胖率进行预测，给出短期预测（到2030年）和长期预测（到2050年）。
  \item 对减肥困难的原因进行分析，对减肥不成功的主要原因进行排序。
  \item 针对几种典型工作生活方式的肥胖人员，提出有针对性的减肥方案。
  \item 在全民体重问题的危机中，对一些典型的危险和机会进行研究。
\end{enumerate}

% ==================== 二、问题分析 ====================
